\section{Discussion}
\label{sec:discussion}
\textbf{Answer to the question.} Recent case counts can flag weeks that will exceed a level
threshold next week with high precision and recall. However, most of this skill comes from
persistence: a simple rule ``we are already above, or close to, the threshold'' is as good as
any model. This is useful for situational awareness but is not early warning. Where early
warning is actually needed (growth and onset surges), skill is modest (growth PR-AUC about
0.26 against a base rate of 0.12) or rests on very few events (S-onset, 5 positives). Weather
from one station adds no measurable value at 1--4 weeks once recent counts are known, which is
consistent with weather acting on dengue through longer lags that are already reflected in the
case trend.

\textbf{Defence points from the brief.}
\emph{Reporting revisions:} two MOH releases differ by at most 8 cases/week, and a simulated
one-week delay reduces PR-AUC only moderately; revisions are therefore not the main risk.
\emph{Immunity and serotypes:} large epidemics such as 2020 and 2022 depend on population
immunity and circulating serotypes, which are not in the data; a model trained on past
epidemics can mistake the timing of a peak, as in the July 2022 failure case.
\emph{Few outbreaks:} onset targets have 5--12 development positives, so their intervals are
wide and one week can change the ranking; we report them as indicative only.
\emph{Cost of false alarms:} at the 0.20 threshold (miss cost = 4 false alarms) the level
targets cost about 10--11 alarms a year and growth about 14; the trade-off for other cost
ratios is in \texttt{alarm\_cost\_tradeoff.csv}.

\textbf{Limitations.} National counts hide local clusters; notified cases depend on testing
and health-seeking behaviour (including the COVID-19 period); there is one weather station and
the humidity series is reanalysis data with only moderate agreement with NEA. The 2022 period
was seen in exploratory runs and the rolling evaluation trains on earlier 2022 outcomes, so
it is a final-period check, not an independent test. All statements are predictive, not
causal.

\textbf{Ethics.} All inputs are aggregated public data, so privacy risk is low. The main risks
are operational: false alarms divert inspection and communication resources, while misses
delay response. The 4:1 cost ratio is our assumption, not elicited from decision-makers. The
scores should support, not replace, human review by public-health staff, and any deployment
would need prospective validation and monitoring for changes in reporting.
